\subsection{All owner paths and even cycles on auxiliary squares}

The owner-size hypothesis in Theorem~\ref{thm:auxiliary-square-parity} can be
removed. We arrange exchanges from different owner cycles simultaneously and
use a forest rank to pay for their combined deficit.

\begin{lemma}\label{lem:square-exchange-forest}
Let $A$ be a union of squares on $V$, and let $B$ have paths and even cycles
as components. Choose at most one edge from each owner cycle whose order is
not divisible by six, with its ends in different auxiliary squares.
Suppose their projections onto the auxiliary squares form a forest with
$r$ edges. There is an $A$-independent $B$-dissociation set of size at least
$\lceil(\diss(B)+r)/2\rceil$.
\end{lemma}
\begin{proof}
For each chosen edge, the prescribed-edge construction in the proof of
Theorem~\ref{thm:auxiliary-square-parity} gives two maximum dissociation sets
of its owner cycle differing exactly at the edge's ends. Keep maximum choices
in all other owner components. The binary choices are independent because
no two chosen edges use the same owner component.

For a resulting maximum owner set $D$, switching one choice toggles the
intersection parities in exactly the two endpoint squares. In a tree on $v$
squares, such edge toggles realize every parity vector with the same total
parity: prescribe the target, and decide the edge at each leaf before removing
that leaf. The remaining root equation holds because the total parity agrees.
Choose all $v$ parities odd when allowed, and otherwise choose $v-1$ odd.
Summing over the trees gives at least $r$ odd square intersections.

In each auxiliary square $K$ choose an independent subset of $D\cap K$ of
size at least $\lceil|D\cap K|/2\rceil$. The union remains dissociation in $B$
by restriction. Its size is at least $(\diss(B)+r)/2$, and integrality gives
the asserted ceiling. Each square's rounding contribution is used once.
\end{proof}

We use the graphic form of the independent representatives theorem of
Rado~\cite{Rado42}, with its proof included. Parallel edges and loops are
allowed. For an edge set $E$, its graphic rank is the largest size of a forest
in $E$, equivalently the number of vertices minus the number of components
in the spanning subgraph with edges $E$.

\begin{lemma}\label{lem:rainbow-forest-rank}
Let $E_1,\ldots,E_c$ be edge families of a finite multigraph. The largest
number $r$ of distinct representatives from different families which form a
forest satisfies
\[
 r=\min_{J\subseteq[c]}
 \left(c-|J|+\operatorname{rank}\left(\bigcup_{i\in J}E_i\right)\right).
\]
\end{lemma}
\begin{proof}
First, full independent representatives exist if and only if
$\operatorname{rank}(E_J)\ge|J|$ for every $J$, where
$E_J=\bigcup_{i\in J}E_i$. Necessity follows by restriction.
For sufficiency induct on $c$; the cases $c=0,1$ are immediate.
If a nonempty proper $J$ is tight, apply induction to its families.
Their representatives form a forest of size $|J|$ spanning $E_J$.
Contract this forest. For any set $K$ of remaining families, their rank in
the contracted graph is
\[
 \operatorname{rank}(E_K\cup E_J)-|J|\ge |K|.
\]
Induction in that graph supplies the other representatives; lifting and
uniting with the contracted forest preserves independence.
If no nonempty proper $J$ is tight, choose a nonloop edge in $E_c$ and
contract it. Every nonempty $K\subseteq[c-1]$ had rank at least $|K|+1$,
and contraction lowers rank by at most one. Induction again applies.

For partial representatives put
$k=\max_J(|J|-\operatorname{rank}(E_J))$, including $J=\varnothing$.
Adjoin $k$ disjoint new edges on new vertices and put all of them in every
family. The full-representative inequalities now hold: every nonempty union
has gained rank $k$. At most $k$ representatives use new edges, leaving at
least $c-k$ original representatives. Conversely every partial system has
at most $c-|J|+\operatorname{rank}(E_J)$ members for each $J$.
These two bounds give the formula.
\end{proof}

\begin{theorem}\label{thm:all-owner-auxiliary-squares}
Let $A$ be a disjoint union of four-cycles on $N=4s$ vertices, and let $B$
be any finite simple graph on the same vertices whose components are paths
or even cycles. There is a set $S$ such that
\[
 A[S]\text{ is edgeless},\qquad \Delta(B[S])\le1,\qquad
 |S|\ge\lceil N/3\rceil.
\]
Consequently every physical partial template whose auxiliary graph is a
union of squares satisfies
\[
 \beta_I(H)\ge\lceil8b/3\rceil.
\]
The coefficient $8/3$ is attained by the thirty-vertex template in
Theorem~\ref{thm:sharp-all-helper}.
\end{theorem}
\begin{proof}
First remove every connected component of $A\cup B$ consisting of one owner
four-cycle on one auxiliary square, or two owner four-cycles on two auxiliary
squares. The owner cycle and auxiliary square in the first type share their four-vertex
support; their edge sets need not coincide. An auxiliary opposite pair gives two compatible
vertices. In the second type each owner cycle crosses the two squares:
a cycle contained in one square would fill it and disconnect that square.
Choose one crossing owner edge. Lemma~\ref{lem:square-exchange-forest}, with
$\diss(B)=4$ and $r=1$, gives three compatible vertices. These components
are closed for both graphs, so their selections can later be united with a
selection on the residual vertices. It suffices to prove the bound on the
residual instance; in what follows $N$ denotes its order.

Call an owner cycle exceptional if its order is not divisible by six.
There are $c$ exceptional cycles. Give a cycle $C_i$ weight
\[
 d_i=\begin{cases}2/3,&|C_i|\equiv4\pmod6,\\
                   1/3,&|C_i|\equiv2\pmod6.\end{cases}
\]
Its order is at least four in the first case and at least eight in the second.
For each $i$, let $E_i$ contain every owner edge of $C_i$ joining different
auxiliary squares, projected onto those squares. Retain physical owner edges
as distinct ground elements, so parallel projected edges remain distinct.
Every residual $E_i$ is nonempty: an owner cycle supported on one square
would be a removed owner four-cycle. The support of one projected owner
cycle is connected, so each label in a subset $J$ lies wholly in one
component of its projected union.

Fix any subset $J$ of exceptional cycles. In a connected component of the
projected union $E_J$, let $x,y$ count the two types and let $v$ count its
auxiliary squares. The owner cycles are disjoint, so their distinct indices
and the capacity of a square give
\[
 4v\ge4x+8y,\qquad v\ge x+2y.
\]
Its graphic rank is $v-1$. If $x+5y\ge3$, then
\[
 v-1\ge x+2y-1\ge(2x+y)/3.
\]
The remaining cases have $y=0$ and $x=1$ or $2$. For $x=1$, $v=1$ would
force a single owner four-cycle filling its auxiliary square, already removed;
otherwise $v-1\ge1\ge2/3$. For $x=2$, $v=2$ would force two owner four-cycles
filling two auxiliary squares. There are no remaining indices there for
another owner component, so this is again a removed closed component;
otherwise $v-1\ge2\ge4/3$.
Thus, for every $J$,
\[
 \operatorname{rank}(E_J)\ge\sum_{i\in J}d_i.
\]
No exceptional cycle has empty support: that would be a coincident owner
four-cycle, one of the removed components.

Lemma~\ref{lem:rainbow-forest-rank} now supplies a forest with at most one
edge from each exceptional owner cycle and size
\[
 r\ge\min_J\left(c-|J|+\sum_{i\in J}d_i\right)
 \ge\sum_{i=1}^c d_i=:d,
\]
since $d_i\le1$. Paths and cycles of order divisible by six have nonnegative
surplus relative to two thirds of their order. The exceptional cycles have
precisely the deficits $d_i$, so
\[
 \diss(B)\ge2N/3-d,\qquad \diss(B)+r\ge2N/3.
\]
Lemma~\ref{lem:square-exchange-forest} gives the residual compatible set of
size at least $\lceil N/3\rceil$. Unite with the two- and three-vertex sets
on the removed four- and eight-vertex components. The sum of these rounded
bounds is at least the ceiling for the total order. Empty instances are
immediate. The physical conclusion follows from
Lemma~\ref{lem:auxiliary-square-normalization}, where $N=2b$; sharpness follows from
Theorem~\ref{thm:sharp-all-helper}.
\end{proof}

For an actual disjoint witness extraction with auxiliary squares, the separate
retained set therefore supplies $Z\ge0$ without an owner-size condition.
The existing ledger gives $51(|S|-\diss(G))\le5|V(G)|$ for that minimal
forming set. As before, a bound on $\SF-\diss$ requires such an extraction
from a maximum-cardinality minimal forming set.
